\metatitle{Transfer matrix method}
\metadescription{The transfer matrix method for finite-state enumeration and rational generating functions.}



\section[transferMatrixMethod]{The transfer matrix method}

The \defin{transfer matrix method} turns a finite-state enumeration problem
into matrix multiplication.  It is especially useful when objects are built
from left to right and only a bounded amount of boundary data is needed to
continue the construction.  See \cite[Sec.~4.7]{StanleyEC1} for the classical
presentation.

Let $\mathcal{S}$ be a finite set of states, and let $A=(a_{st})_{s,t\in\mathcal{S}}$
be a matrix with entries in a commutative ring $R$.  We interpret $a_{st}$ as
the weight of a transition from state $s$ to state $t$.

\begin{proposition}[Transfer matrix method]
Let $\alpha$ and $\beta$ be column vectors indexed by $\mathcal{S}$, and define
\[
  f_n \coloneqq \alpha^\top A^n \beta .
\]
Then the ordinary generating function of the sequence $(f_n)_{n\geq 0}$ is
\[
  \sum_{n\geq 0} f_n z^n
  =
  \alpha^\top (I-zA)^{-1}\beta .
\]
In particular, if $R$ is a field, then this generating function is rational.
\end{proposition}

\begin{proof}
The identity follows by summing the geometric series
\[
  I+zA+z^2A^2+\dotsb=(I-zA)^{-1}.
\]
The claim follows.
\end{proof}

The method also applies with extra variables in the entries of $A$.  In that
case the matrix records a weighted enumeration, and specializations of the
variables recover ordinary counts.


\subsection[transferMatrixFibonacciWords]{Binary words with no consecutive ones}

Let $a_n$ be the number of binary words of length $n$ with no consecutive ones.
We use two states, recording whether the last letter is $0$ or $1$.  If
\[
  b_n =
  \begin{pmatrix}
  b_n^{(0)}\\
  b_n^{(1)}
  \end{pmatrix}
\]
records the number of valid words ending in $0$ and $1$, then
\[
  b_{n+1}
  =
  \begin{pmatrix}
  1 & 1\\
  1 & 0
  \end{pmatrix}
  b_n .
\]
Thus $a_n=b_n^{(0)}+b_n^{(1)}$, and one obtains
\[
  \sum_{n\geq 0} a_n z^n
  =
  \frac{1}{1-z-z^2}.
\]
Equivalently, $a_n=F_{n+2}$, where $F_n$ is the Fibonacci sequence.


\subsection[transferMatrixGraphColorings]{Proper colorings of paths and cycles}

Let $B$ be the $q\times q$ matrix with $0$ on the diagonal and $1$ off the
diagonal.  A product of entries in $B$ records whether adjacent vertices have
distinct colors.

For the path graph on $n$ vertices, the number of proper $q$-colorings is
\[
  \mathbf{1}^\top B^{n-1}\mathbf{1}
  =
  q(q-1)^{n-1}.
\]
For the cycle graph $C_n$, the first and last colors must also be adjacent,
so the transfer matrix expression is a trace:
\[
  \chi_{C_n}(q)
  =
  \operatorname{tr}(B^n)
  =
  (q-1)^n + (q-1)(-1)^n .
\]
This is the usual chromatic polynomial of a cycle.


\subsection[transferMatrixDominoTilings]{Domino tilings of a strip}

The number of domino tilings of a $2\times n$ rectangle can be computed by
recording which cells in the next column are already occupied by horizontal
dominoes.  The possible boundary states are the four subsets of the two cells,
which we encode by $0,1,2,3$.  With rows and columns ordered in this way, the
transfer matrix is
\[
  A=
  \begin{pmatrix}
  1 & 0 & 0 & 1\\
  0 & 0 & 1 & 0\\
  0 & 1 & 0 & 0\\
  1 & 0 & 0 & 0
  \end{pmatrix}.
\]
The entry $(A^n)_{0,0}$ counts tilings which start and end with an empty
boundary.  This gives the Fibonacci recurrence for the number of domino
tilings of a $2\times n$ rectangle.

Weighted versions of the same matrix give bivariate generating functions
recording, for example, the number of vertical and horizontal dominoes.


\subsection[transferMatrixWeightedTilings]{Weighted tilings by larger tiles}

Let $P_n(t)$ be the generating polynomial for tilings of a $4\times n$
rectangle using unit squares and $2\times 2$ squares, where each $2\times 2$
square contributes a factor of $t$. Thus the coefficient of $t^k$ counts
tilings using exactly $k$ large squares.
This is the rectangular big-tile case $\mu=(2,2)$ in the tiling polynomials
studied by \name{John Ahlberg} and \name{Per Alexandersson}
\cite{AhlbergAlexandersson2026}.

Process the board column by column. A boundary state records which cells in
the next column have already been occupied by $2\times 2$ squares started in
the previous column. For height $4$, starting from the empty boundary, only
the five states
\[
  \emptyset,\quad 12,\quad 23,\quad 34,\quad 1234
\]
can occur, where for instance $23$ means that rows $2$ and $3$ are already
occupied.

Using column vectors, with rows indexed by the next state and columns indexed
by the current state in the order above, the transfer matrix is
\[
  A(t)=
  \begin{pmatrix}
  1 & 1 & 1 & 1 & 1 \\
  t & 0 & 0 & t & 0 \\
  t & 0 & 0 & 0 & 0 \\
  t & t & 0 & 0 & 0 \\
  t^2 & 0 & 0 & 0 & 0
  \end{pmatrix}.
\]
For example, from the empty state one may start no large square, one large
square in rows $12$, $23$, or $34$, or two large squares in rows $12$ and
$34$. This explains the first column of $A(t)$.

The desired polynomial is the empty-to-empty entry
\[
  P_n(t)=\bigl(A(t)^n\bigr)_{\emptyset,\emptyset}.
\]
Consequently,
\[
  \sum_{n\geq 0} P_n(t) z^n
  =
  \frac{1-tz}
  {1-(1+t)z-(2t+t^2)z^2+(t^2+t^3)z^3}.
\]
The first few polynomials are
\[
\begin{array}{c|l}
n & P_n(t) \\ \hline
0 & 1\\
1 & 1\\
2 & 1+3t+t^2\\
3 & 1+6t+4t^2\\
4 & 1+9t+16t^2+8t^3+t^4\\
5 & 1+12t+37t^2+34t^3+9t^4 .
\end{array}
\]
